\documentclass[10pt,a4paper]{amsart}
\usepackage[margin=19mm]{geometry}
\usepackage{amsmath,amssymb,mathtools}
\usepackage[scr=rsfs,cal=euler]{mathalfa}
\usepackage{xfrac}
\usepackage[unicode,hidelinks,bookmarks=false]{hyperref}
\newcommand{\ie}{i.e.\ }
\newcommand{\cf}{cf.\ }
\newcommand{\resp}{respectively\ }
\newcommand{\Subsection}{\subsection}
\providecommand{\nobreakdash}{\nobreak\hbox{-}}
\setlength{\emergencystretch}{2em}
\multlinegap0pt
\usepackage{amssymb}
\usepackage[all]{xy}
\newtheorem{lem}{Lemma}
\newcommand{\Z}{\mathbb{Z}}
\DeclareMathOperator{\lk}{lk}
\title{Non-isotopic transverse tori in Engel manifolds}
\author{Marc Kegel}
\date{Source: arXiv:2205.04853v2 (CC BY 4.0)}
\begin{document}
\maketitle

\noindent\emph{Source and licence.} Non-isotopic transverse tori in Engel manifolds. Version arXiv:2205.04853v2 (CC BY 4.0), distributed by arXiv under the Creative Commons Attribution 4.0 International licence, http://creativecommons.org/licenses/by/4.0/, as recorded in the arXiv licence metadata for this exact version and shown on the abstract page https://arxiv.org/abs/2205.04853v2. Full licence notice: \texttt{licenses/arxiv-2205.04853-CC-BY-4.0.txt}.\par\smallskip

\noindent\textit{Editorial scope.} Counted unit: the article's lem lem:hom2 (source0.tex lines 400--402). The article's complete original proof of it is delivered with it (source0.tex lines 404--423, one \texttt{proof} environment of 20 lines). No mathematics is written, repaired or re-derived: the statement and the proof are the article's own lines, in the article's own notation, and no retained line is substituted or re-flowed.  No further source-local context is needed: every cross-reference of every retained line resolves inside the two delivered spans.  Nothing was omitted from any retained span, so the package carries no omission record.  No retained line contains a citation, so the wrapper carries no bibliography and no bibliography entry.  No retained line contains a figure environment, an image-inclusion command or drawing code, so no image is delivered and the package declares no asset dependency.  Editorial formatting only, in addition: an AMS \texttt{amsart} preamble that restates the article's own declarations and macros used by the retained lines, namely the environments \texttt{lem} on separate counters, together with the standard packages those lines need.  The retained environments print in this document as Lemma 1 in order of appearance, whereas the article prints the same items with the numbering of its own full document.  The complete native source of the article as distributed in the arXiv e-print for this exact version is bundled unchanged as \texttt{sources/125525/source0.tex}.
\begin{lem} \label{lem:hom2}
Let $K$ and $K'$ be nullhomologous knots in a closed $3$-manifold $N$. We consider $L=S^1\times K$ and $L'=S^1\times K'$ in $M=S^1\times N$. Then $H_2(M\setminus L)$ is isomorphic to $H_1(N)\oplus H_2(N)\oplus \Z$ and the linking class $\lk(L',L)$ takes values in the last free abelian group generated by $S^1\times \mu_K$ where $\mu_K$ is the meridian of $K$.
\end{lem}

\begin{proof}
First we combine Poincar\'e duality, excision and the universal coefficient theorem to deduce that
\begin{equation*}
H_2(M\setminus \mathring{\nu L})\cong F_2(M,L)\oplus T_1(M,L),
\end{equation*}
where $F_2(M,L)$ denotes the free part of $H_2(M,L)$ and $T_1(M,L)$ the torsion of $H_1(M,L)$. By applying the relative K\"unneth formula to $S^1\times (N,K)$ we compute
\begin{equation*}
H_k(M,L)\cong H_{k-1}(N,K)\oplus H_k(N,K).
\end{equation*}
Next, we consider the long exact sequence of the pair $(N,K)$
\begin{equation*}
0\longrightarrow H_2(N)\longrightarrow H_2(N,K)\longrightarrow \Z\overset{0}{\longrightarrow} H_1(N)\longrightarrow H_1(N,K)\longrightarrow 0,
\end{equation*}
where the map $\Z\rightarrow H_1(N)$ is the $0$-map since $K$ is nullhomologous in $N$. Thus 
\begin{equation*}
H_1(N,K)=H_1(N)\,\text{ and }\, H_2(N,K)\cong H_2(N)\oplus\Z.
\end{equation*}
Putting everything together and using that $H_2(N)$ is a free group we obtain that $H_2(M\setminus L)$ is isomorphic to $H_1(N)\oplus H_2(N)\oplus \Z$, where we can follow the above isomorphisms to see that the $\Z$-summand is generated by $S^1\times \mu_K$.
Since $L'$ is nullhomologous in $H_2(M)$ and $K'$ is nullhomologous in $H_1(N)$ it follows that $L'$ seen as a homology class in $H_2(M\setminus L)$ is a multiple of $S^1\times \mu_K$.
\end{proof}

\end{document}
